% What crosses the trust boundary, as a table shared by chapters 5 and 6 (one source).
\begin{table}[H]
\centering
\caption{What crosses the trust boundary. Bob can check every row from outside Alice's harbor except the last: a signed, intact history does not prove who controls the key. The market needs that row and does not have it yet. [internal]}
\label{tab:keystone-split}
\small
\renewcommand{\arraystretch}{1.15}
\begin{tabularx}{\linewidth}{@{}>{\raggedright\arraybackslash}X >{\raggedright\arraybackslash}X >{\raggedright\arraybackslash}X@{}}
\toprule
\textbf{Claim} & \textbf{Inside Alice's trusted daemon} & \textbf{From Bob's foreign view} \\
\midrule
the actor key was minted by the daemon & known: the daemon minted it & the key remains visible; the signature checks \\
the history is append-only & known: the daemon wrote it & log integrity remains verifiable; inclusion proofs check \\
the registry is local & known & the foreign registry is inspectable, not trusted \\
the key is bound to an accountable principal & known: the daemon's own binding & \textbf{no cross-operator binding proof}: nothing Bob can check says who controls the key \\
\bottomrule
\end{tabularx}
\end{table}
